% test-028.tex -- comandos \spcase y \spOmega
%
% Compilar:  lualatex-dev test-028.tex
% Validar:   show-pdf-tags --xml test-028.pdf | rnv-wrapp
%            verapdf --flavour ua2 --format text test-028.pdf
%
% Cada caso es un parrafo numerado, para poder referirse a el al
% escuchar el PDF. Los comentarios "doc:" indican la lectura que
% documenta el manual de usuario; en el resto, la lectura se revisa
% escuchando. Las entradas invalidas (que detienen la compilacion)
% no van aqui.
\DocumentMetadata
  {
    lang=es-CL,
    pdfversion=2.0,
    pdfstandard=ua-2,
    tagging=on,
    tagging-setup={math/setup={mathml-SE}},
  }
\documentclass{article}
\usepackage[spanish]{babel}
\usepackage{unicode-math}
\usepackage{spintent}
\usepackage{hyperref}
\hypersetup
  {
    pdfdisplaydoctitle = true,
    pdftitle           = {test-028: probabilidad},
  }
\pagestyle{empty}
\begin{document}

\section*{A. El comando spcase}

% doc: «casos favorables»
1. Favorables: $\spcase{f}$.

% doc: «casos posibles»
2. Posibles: $\spcase{p}$.

% doc: «casos totales»
3. Totales: $\spcase{t}$.

% doc: «casos que cumplen»
4. Con read-arg: $\spcase[read-arg={casos que cumplen}]{f}$.

% doc: sin lectura
5. Con silent: $\spcase[silent=true]{p}$.

\section*{B. El comando spOmega}

% doc: «espacio muestral»
6. Espacio muestral: $\spOmega$.

% doc: «cardinalidad del espacio muestral»
7. Cardinalidad: $\spOmega[card=true]$.

% doc: «espacio muestral sub 1»
8. Con sub-index: $\spOmega[sub-index={1}]$.

% doc: «cardinalidad del espacio muestral sub 1»
9. Cardinalidad con sub-index: $\spOmega[card=true,sub-index={1}]$.

% doc: «espacio muestral de 1»
10. Con read-sub=false: $\spOmega[sub-index={1},read-sub=false]$.

% doc: «universo»
11. Con read-arg: $\spOmega[read-arg={universo}]$.

% doc: sin lectura
12. Con silent: $\spOmega[silent=true]$.

\section*{C. Dentro de fórmulas}

% doc: «p es igual a la fracción con numerador casos favorables y denominador casos posibles»
13. Regla de Laplace: $P = \dfrac{\spcase{f}}{\spcase{p}}$.

% doc: «p es igual a la fracción con numerador casos favorables y denominador cardinalidad del espacio muestral»
14. Con la cardinalidad: $P = \dfrac{\spcase{f}}{\spOmega[card=true]}$.

% doc: «casos totales es igual a cardinalidad del espacio muestral»
15. Casos totales: $\spcase{t} = \spOmega[card=true]$.

\end{document}
